\documentclass[11pt,a4paper]{article}
\usepackage[T1]{fontenc}
\usepackage[utf8]{inputenc}
\usepackage{lmodern,amsmath,amssymb}
\usepackage[margin=25mm]{geometry}
\usepackage[hidelinks]{hyperref}
\hypersetup{pdftitle={An Agmon bound for the Kogut--Susskind ground state: large fields are suppressed at rate $},pdfauthor={navstokgap research working notes}}
\newcommand{\tightlist}{\setlength{\itemsep}{0pt}\setlength{\parskip}{0pt}}
\setlength{\emergencystretch}{3em}
\setcounter{secnumdepth}{0}
\begin{document}
\hypertarget{an-agmon-bound-for-the-kogutsusskind-ground-state-large-fields-are-suppressed-at-rate-1g2-per-plaquette}{%
\section{\texorpdfstring{An Agmon bound for the Kogut--Susskind ground
state: large fields are suppressed at rate \(1/g^2\) per
plaquette}{An Agmon bound for the Kogut--Susskind ground state: large fields are suppressed at rate 1/g\^{}2 per plaquette}}\label{an-agmon-bound-for-the-kogutsusskind-ground-state-large-fields-are-suppressed-at-rate-1g2-per-plaquette}}

The estimate that \href{flow-before-decimation.md}{the
flow-before-decimation note} identified as the missing ingredient exists
and is standard: Agmon's method applies to the Kogut--Susskind
Hamiltonian because its configuration space is a compact Riemannian
manifold, its magnetic term is a bounded potential and its ground state
is positive. Writing
\[H=\frac{\hbar c}{a}\Big[\underbrace{\frac{g^2}{2}}_{A'}\,(-\Delta)
+\underbrace{\frac{2}{g^2}}_{B'}\,V\Big],\qquad
V(U)=\sum_p\big(N-\operatorname{Re}\operatorname{tr}U_p\big),\] the
identity \[A'\!\int\!\big|\nabla(e^{\rho}\Omega)\big|^2
+\int\!\big(B'V-e_0-A'|\nabla\rho|^2\big)e^{2\rho}\Omega^2=0,
\qquad e_0=\frac{aE_0}{\hbar c},\] holds for every Lipschitz \(\rho\)
and gives exponential decay of the ground state into the region where
the magnetic energy exceeds its mean. The Agmon metric is
\(\sqrt{(B'V-e_0)/A'}\;|dU|\), with \(B'/A'=4/g^4\), and a configuration
carrying an excess of \(n\) excited plaquettes sits at Agmon distance of
order \(2n\sqrt v/g^2\) from the allowed region, where \(v\) is the
excess per plaquette. Hence
\[\big|\Omega(U)\big|^2\ \lesssim\ \exp\Big[-\frac{c\,n\sqrt N}{g^2}\Big],\]
the ground-state measure suppressing large-field regions at a rate
\(\propto1/g^2\) per plaquette, \textbf{the same coupling dependence as
the Euclidean Wilson weight} \(e^{-S_w}\) with \(\beta=2N/g^2\). The
suppression is strong exactly at weak coupling, which is where the
renormalization steps of the route live, and weak at strong coupling,
where the cluster expansion already works. This supplies the Hamiltonian
counterpart of the large-field estimate whose absence
\href{typical-field-strength-window.md}{the typical-field note}
recorded, and it does so by a standard technique rather than a new one.
Constants explicit; the distance estimate is a scaling argument and is
labelled as such; nothing promoted.

\begin{quote}
\textbf{Corrected.} The distance estimate of Section 2 is carried out
with constants in \href{agmon-global-not-local.md}{the global/local
note}, and it controls the \textbf{global} excess of magnetic energy,
not a local region: the gradient norm
\(\|\nabla V\|_\infty\le2\sqrt{2N|\mathcal E|}\) is extensive, so a
fixed local excess in a large volume gives a bound that degrades like
\(|\mathcal P|^{-1/2}\). Read ``\(n\) excess plaquettes'' below as a
global excess. The identity of Lemma 1 and the estimate of Corollary 2
are unaffected.
\end{quote}

\hypertarget{setting-and-the-ground-state-identity}{%
\subsection{1. Setting and the ground-state
identity}\label{setting-and-the-ground-state-identity}}

The configuration space \(\mathcal M=G^{\mathcal E}\) is a compact
Riemannian manifold with the product bi-invariant metric, of dimension
\(|\mathcal E|\dim G\), and \(-\Delta\) is its Laplace--Beltrami
operator. By T1 of \href{mass-gap-obligations-lattice.md}{the
obligations map} the ground state \(\Omega\) is strictly positive and
satisfies \[-A'\Delta\Omega+\big(B'V-e_0\big)\Omega=0,\qquad
A'=\frac{g^2}{2},\quad B'=\frac{2}{g^2},\quad e_0=\frac{aE_0}{\hbar c},\]
after dividing by \(\hbar c/a\).

\textbf{Lemma 1 (Agmon identity).} For every Lipschitz
\(\rho:\mathcal M\to\mathbb R\),
\[A'\!\int_{\mathcal M}\big|\nabla(e^{\rho}\Omega)\big|^2\,d\mu
+\int_{\mathcal M}\big(B'V-e_0-A'|\nabla\rho|^2\big)\,e^{2\rho}\,\Omega^2\,d\mu=0 .\]

\emph{Proof.} Multiply the eigenvalue equation by \(e^{2\rho}\Omega\)
and integrate; there is no boundary term because \(\mathcal M\) is
closed. Using
\(\int(-\Delta\Omega)e^{2\rho}\Omega=\int\nabla\Omega\cdot\nabla(e^{2\rho}\Omega)\)
and the pointwise identity
\[\big|\nabla(e^\rho\Omega)\big|^2=e^{2\rho}\big|\nabla\Omega\big|^2
+2e^{2\rho}\Omega\,\nabla\rho\cdot\nabla\Omega+e^{2\rho}\Omega^2|\nabla\rho|^2
=\nabla\Omega\cdot\nabla(e^{2\rho}\Omega)+e^{2\rho}\Omega^2|\nabla\rho|^2,\]
the first term becomes
\(\int|\nabla(e^\rho\Omega)|^2-\int e^{2\rho}\Omega^2|\nabla\rho|^2\).
Adding \(\int(B'V-e_0)e^{2\rho}\Omega^2\) gives the identity.
\(\square\)

\textbf{Corollary 2 (Agmon estimate).} Let
\(\Sigma_\epsilon=\{U:B'V(U)-e_0\ge\epsilon\}\) and let \(\rho\) be
Lipschitz with \(A'|\nabla\rho|^2\le B'V-e_0-\epsilon/2\) on
\(\Sigma_\epsilon\) and \(\rho\le\rho_{\max}\) on
\(\mathcal M\setminus\Sigma_\epsilon\). Then
\[\frac\epsilon2\int_{\Sigma_\epsilon}e^{2\rho}\,\Omega^2\,d\mu
\ \le\ \big(e_0+A'\|\nabla\rho\|_\infty^2\big)\,e^{2\rho_{\max}}\!\!\int_{\mathcal M\setminus\Sigma_\epsilon}\!\!\Omega^2\,d\mu
\ \le\ C\,e^{2\rho_{\max}} .\]

\emph{Proof.} Drop the nonnegative gradient term in Lemma 1 and split
the integral over \(\Sigma_\epsilon\) and its complement, where the
integrand is bounded below by \(-(e_0+A'\|\nabla\rho\|_\infty^2)\).
\(\square\)

The admissible weights are exactly the functions with
\(|\nabla\rho|\le\sqrt{(B'V-e_0)/A'}\), so the largest is the
\textbf{Agmon distance}
\[d(U)=\inf_{\gamma:\,U\to\{B'V\le e_0\}}\int_\gamma\sqrt{\frac{B'V-e_0}{A'}}\;|d\gamma| ,\]
and Corollary 2 with \(\rho=(1-\delta)d\) gives
\(\int e^{2(1-\delta)d}\Omega^2\le C(\delta)\).

\hypertarget{the-agmon-distance-of-a-large-field-region}{%
\subsection{2. The Agmon distance of a large-field
region}\label{the-agmon-distance-of-a-large-field-region}}

Write \(\bar v=e_0/B'\) for the mean magnetic energy, which satisfies
\(\bar v\le\langle V\rangle_{\rm Haar}=N|\mathcal P|\) because
\(E_0\le\langle1,H1\rangle\) and
\(\langle\operatorname{Re}\operatorname{tr}U_p\rangle_{\rm Haar}=0\).
Consider a configuration \(U\) whose magnetic energy exceeds \(\bar v\)
by an amount carried on \(n\) plaquettes, each in excess by \(v=O(N)\),
so \(V(U)-\bar v\simeq nv\).

\emph{Flat length.} Relaxing those \(n\) plaquettes requires moving of
order \(n\) link variables by an angle of order one. In the product
metric the flat distance is of order \(\sqrt n\), since \(n\)
coordinates each move \(O(1)\).

\emph{Agmon weight.} Along such a path the excess decreases from \(nv\)
to \(0\), with mean of order \(nv/2\), so the weight
\(\sqrt{(B'V-e_0)/A'}=\sqrt{(B'/A')(V-\bar v)}\) is of order
\(\sqrt{(4/g^4)\,nv/2}=\sqrt{2nv}/g^2\).

\emph{Distance.} Multiplying, and keeping only the order,
\[d(U)\ \simeq\ \frac{\sqrt{2nv}}{g^2}\cdot\sqrt n\ =\ \frac{n\sqrt{2v}}{g^2}.\]

\textbf{Proposition 3 (scaling form).} With the distance estimate above,
Corollary 2 gives, for the ground-state measure,
\[\int_{\{n\ \rm excess\ plaquettes\}}\Omega^2\,d\mu
\ \lesssim\ C\exp\Big[-\frac{2(1-\delta)\sqrt{2v}}{g^2}\,n\Big],\]
exponential in the number of excited plaquettes with rate
\(\lambda=2\sqrt{2v}/g^2\) per plaquette, \(v=O(N)\).

The flat-length and mean-excess steps are scaling estimates; the
identity of Lemma 1 and the estimate of Corollary 2 are exact.

\hypertarget{comparison-with-the-euclidean-weight}{%
\subsection{3. Comparison with the Euclidean
weight}\label{comparison-with-the-euclidean-weight}}

The Euclidean Wilson measure is \(e^{-S_w}\,d\mu\) with
\[S_w=\frac{2N}{g_E^2}\sum_p\Big(1-\frac1N\operatorname{Re}\operatorname{tr}U_p\Big)
=\frac{2}{g_E^2}\,V,\] suppressing a configuration with \(n\) excited
plaquettes by \(e^{-2nv/g_E^2}\). Proposition 3 gives
\(e^{-2n\sqrt{2v}/g^2}\) for the Hamiltonian ground-state measure. The
two agree in their dependence on the coupling, \(1/g^2\) per excited
plaquette, and differ in the power of the excess per plaquette, linear
against square-root, which is the usual difference between a Boltzmann
weight and a WKB weight: the Agmon bound is a tunnelling estimate and
the Wilson weight is a direct cost.

\textbf{The direction is the useful one.} The suppression rate grows as
\(g\to0\), so the large-field regions that
\href{typical-field-strength-window.md}{the typical-field note} could
not control in the Hamiltonian framework are exponentially rare in the
ground-state measure precisely at weak coupling, which is where the
renormalization steps of \href{mass-gap-position.md}{the position note}
§6 are needed. At strong coupling the rate degrades, and there the
cluster expansion of \href{strong-coupling-uniform-gap.md}{the T2 note}
already applies.

\hypertarget{what-this-supplies-and-what-it-still-lacks}{%
\subsection{4. What this supplies and what it still
lacks}\label{what-this-supplies-and-what-it-still-lacks}}

\emph{Supplies.} A Hamiltonian counterpart of the large-field estimate,
by a standard technique, with an explicit rate. It converts the
state-relative criterion demanded by
\href{flow-before-decimation.md}{the flow-before-decimation note} §3
into a computable statement: expectations of quantities that grow
exponentially in the local field strength are finite in the ground state
as long as their growth rate stays below \(\lambda=2\sqrt{2v}/g^2\).

\emph{Lacks.} Three things. The estimate is for the ground state, and
the route needs it for the whole low-energy subspace, which requires
either the same argument for low-lying eigenfunctions, available with
\(e_0\) replaced by \(e_1\) at the cost of a smaller forbidden region,
or a spectral-projection version. The distance estimate of Section 2 is
a scaling argument and needs a proof with constants, which is a
geometric computation on \(G^{\mathcal E}\). And the resulting
suppression must be matched against the growth \(e^{2t\|G\|_\infty}\) of
the flow Jacobian: by \href{lattice-truncation-uniform.md}{the
lattice-truncation note} that factor is a pure number per doubling, so
the match is comfortable within one step and the question is again the
composition of steps.

\hypertarget{consequence-for-state}{%
\subsection{5. Consequence for STATE}\label{consequence-for-state}}

The missing ingredient of the previous note has a standard source. The
Kogut--Susskind ground state obeys an Agmon bound whose rate,
\(\lambda\simeq2\sqrt{2v}/g^2\) per excess plaquette, reproduces the
coupling dependence of the Euclidean Wilson weight and is strong exactly
at weak coupling. The next items, in order of tractability: prove the
distance estimate of Section 2 with constants; extend the bound from the
ground state to the spectral subspace below the gap; then assemble the
state-relative decimation criterion that these two make possible, which
is the first version of the renormalization step in which the difficulty
is located in the composition of steps rather than in any single one.
\end{document}
